\section{Truncated Hermite hierarchies}\label{sec:truncated}

Fourier--Hermite codes for KRMHD evolve the Hermite moments $g_0,\dots,g_M$ of a compressive branch and close the hierarchy at $m=M$; two examples are Viriato \citep{Loureiro2016} and GANDALF \citep{Kanekar2025}. This section shows that the truncated hierarchy keeps the slice structure of section~\ref{sec:slices} exactly. It consists of $M+1$ passive slices, and with the zero closure their phase velocities and weights are the nodes and weights of the Gauss quadrature rule of the spectral density $\varrho\,\dd u$ of (\ref{eq:varrho}). We then derive the conserved quadratic forms, the discrete phase-space helicity at odd and even $M$, the distribution of the phase velocities and the largest of them, the coupling of the slices by hypercollisions, the reflection of phase-mixing packets at the closure, and the consequences for explicit time stepping. Appendix~\ref{app:truncated} gives the proofs and further detail, and appendix~\ref{app:timestepping} analyses GANDALF's time steps in slice variables. Unless stated otherwise, every exact identity of this section was checked to round-off in double precision, against GANDALF's right-hand side, time step or matrices (version 0.6.0) where these are involved, and otherwise symbolically or by direct quadrature (appendix~\ref{app:truncated}); the asymptotic laws were checked numerically as stated. Section~\ref{sec:numerics} shows the main checks.

\subsection{The slices of the truncated hierarchy}\label{tr:slicesec}

For one branch truncated at $m=M$, the vector $g=(g_0,\dots,g_M)^{\mathsf T}$ obeys
\begin{equation}\label{tr:hierarchy}
  \partial_tg+\{\Phi,g\}+s\,\gradpar(Ag)=0,\qquad \gradpar=\partial_z+\{\Psi,\cdot\},
\end{equation}
where $A$ is an $(M+1)\times(M+1)$ ladder matrix. With the zero closure, $g_{M+1}=0$, $A$ consists of the first $M+1$ rows and columns of the ladder matrix of section~\ref{sec:model}: it is tridiagonal with zero diagonal, and $A_{m,m+1}=A_{m+1,m}=\sqrt{(m+1)/2}$ except $A_{10}=(1-1/\Lambda)/\sqrt2$. GANDALF also offers the closure $g_{M+1}=g_{M-1}$, which it calls symmetric and which we call the copy closure. For $M\ge2$ it changes the single entry $A_{M,M-1}$ to $\sqrt{M/2}+\sqrt{(M+1)/2}$. GANDALF applies $\partial_z$ and $\{\Psi,\cdot\}$ to the same combination of neighbouring moments, so the closure enters the streaming and the field-line term in the same way. The weight $\Pw=\mathrm{diag}(p_m)$, with $p_0=1-1/\Lambda$ and $p_m=1$ for $m\ge1$, makes $\Pw A$ symmetric for the zero closure. For the copy closure $\Pw A$ is symmetric when the last weight is replaced by
\begin{equation}\label{tr:pM}
  p_M=\frac{\sqrt M}{\sqrt M+\sqrt{M+1}}\approx\frac12 ,
\end{equation}
and $\Pw A$ is then the same symmetric matrix for both closures (appendix~\ref{app:truncated}). In this section $\Pw$ denotes the weight of whichever closure is used, and $D_c=\mathrm{diag}(1,\dots,1,p_M)$ is the factor by which the copy weight differs from the zero-closure weight. The weight is positive definite if and only if $\Lambda\notin[0,1]$.

In GANDALF each bracket is evaluated pseudo-spectrally and dealiased by the 2/3 rule, the right-hand side is multiplied by the dealiasing mask, and its $\boldsymbol k=0$ component is set to zero. Each of these operations acts on every $g_m$ in the same way. GANDALF evolves the Elsasser potentials $z^\pm$ together with $g$, but the statements below hold for any $\Phi$ and $\Psi$.

\begin{theorem}\label{thm:trunc}
Let $\Lambda\notin[0,1]$, the value $\Lambda=\infty$ included, and let the hierarchy be truncated at $m=M$ with the zero closure or the copy closure.
\begin{itemize}
\item[(i)] $A$ has $M+1$ real and simple eigenvalues $u_0<u_1<\dots<u_M$. The right eigenvector of $u_j$ is $\boldsymbol r_j=(r_0(u_j),\dots,r_M(u_j))^{\mathsf T}$, where $r_m(u)$, the $m$th Hermite coefficient of the Case--Van Kampen mode $F_u$ of section~\ref{sec:slices}, is a polynomial of degree $m$. The left eigenvector is $\boldsymbol w_j=\Pw\boldsymbol r_j$, and $\boldsymbol w_j^{\mathsf T}\boldsymbol r_l=n_j\,\delta_{jl}$ with $n_j=\boldsymbol r_j^{\mathsf T}\Pw\boldsymbol r_j>0$.
\item[(ii)] Let $\Phi$ and $\Psi$ be arbitrary real fields on the grid, and let the brackets in (\ref{tr:hierarchy}) be evaluated by any bilinear rule, possibly followed by linear operations, such as dealiasing, that act on every $g_m$ in the same way. Define the unit-weight slice amplitudes
\begin{equation*}
  a_j=\frac{\boldsymbol w_j^{\mathsf T}g}{\sqrt{n_j}}=\sum_{m=0}^Mp_m\,R_{mj}\,g_m,\qquad R_{mj}=\frac{r_m(u_j)}{\sqrt{n_j}} .
\end{equation*}
Then $g$ solves (\ref{tr:hierarchy}) if and only if, for $j=0,\dots,M$,
\begin{equation}\label{eq:truncslice}
  \partial_ta_j+\{\Phi+su_j\Psi,\,a_j\}+su_j\,\partial_za_j=0,
\end{equation}
with the brackets evaluated by the same rule.
\item[(iii)] $R^{\mathsf T}\Pw R=I$ and $g_m=\sum_jR_{mj}\,a_j$. For every Fourier mode, $\sum_mp_m|g_m(\boldsymbol k)|^2=\sum_j|a_j(\boldsymbol k)|^2$.
\end{itemize}
\end{theorem}

\begin{proof}
(i) Since $\Pw A$ is symmetric and $\Pw$ is positive definite, $\Pw^{1/2}A\,\Pw^{-1/2}$ is a real symmetric tridiagonal matrix. Its off-diagonal entries $(A_{m,m+1}A_{m+1,m})^{1/2}$ are non-zero, because $A_{01}A_{10}=(1-1/\Lambda)/2>0$ and the other products are positive for both closures. A symmetric tridiagonal matrix with non-zero off-diagonal entries has real and simple eigenvalues, and so has $A$, which is similar to it. Rows $0$ to $M-1$ of $A\boldsymbol r=u\boldsymbol r$ are the three-term recurrence of the polynomials $r_m$ (appendix~\ref{app:truncated}), which fixes the eigenvector up to a factor, chosen so that its first component is $r_0=1$; row $M$ is the condition on the eigenvalue. Since $\Pw A=A^{\mathsf T}\Pw$, $\boldsymbol w_j^{\mathsf T}A=\boldsymbol r_j^{\mathsf T}\Pw A=u_j\boldsymbol w_j^{\mathsf T}$. The vectors $\Pw^{1/2}\boldsymbol r_j$ are eigenvectors of a symmetric matrix with distinct eigenvalues, so $\boldsymbol r_j^{\mathsf T}\Pw\boldsymbol r_l=0$ for $j\neq l$, and $n_j>0$ because $\Pw$ is positive definite.

(ii) The amplitude $a_j$ is a fixed linear combination of the $g_m$. The brackets with $\Phi$ and with $\Psi$, the derivative $\partial_z$ and the subsequent linear operations act on every $g_m$ in the same way, so they commute with this combination, and $\boldsymbol w_j^{\mathsf T}A=u_j\boldsymbol w_j^{\mathsf T}$ gives $\sum_mp_mR_{mj}(Ag)_m=u_ja_j$. Applying the combination to (\ref{tr:hierarchy}) gives $\partial_ta_j+\{\Phi,a_j\}+su_j\,\partial_za_j+s\{\Psi,u_ja_j\}=0$, which is (\ref{eq:truncslice}). Conversely, by (iii) each $g_m$ is a fixed linear combination of the $a_j$, and the same argument returns (\ref{tr:hierarchy}).

(iii) By (i) and the normalisation of $R$, $R^{\mathsf T}\Pw R=I$, so $R^{-1}=R^{\mathsf T}\Pw$ and $g=Ra$. Then $\sum_mp_m|g_m|^2=a^\dagger R^{\mathsf T}\Pw R\,a=\sum_j|a_j|^2$.
\end{proof}

Each slice $a_j$ is a passive scalar carried by the divergence-free velocity $\zhat\times\nabla_\perp(\Phi+su_j\Psi)+su_j\zhat=\uE+su_j\bhat$, which is the velocity $\boldsymbol U_u$ of (\ref{eq:Uu}) at $u=u_j$. Theorem~\ref{thm:trunc} is therefore the truncated form of theorem~\ref{thm:slice}, with the continuum of phase velocities replaced by the $M+1$ eigenvalues of $A$. It holds for any $z^\pm$, whether or not they obey the RMHD equations, for both of GANDALF's closures and for every $\Lambda\notin[0,1]$. Its proof uses only the bilinearity of the discrete bracket, not the dealiasing. We checked (\ref{eq:truncslice}) against GANDALF's right-hand side for $M$ between 2 and 65, $\Lambda=\pm\sqrt5$ and $3$, $\beta_i$ between $0.37$ and $1.7$, and both closures, in states in which the brackets supply up to the whole of the right-hand side. The relative mismatch is at most $2\times10^{-13}$, and it is $1.6\times10^{-14}$ when $z^\pm$ carry power outside the dealiasing band. Advecting each slice by $\Phi$ alone, without $su_j\Psi$, gives mismatches of $0.6$ to $0.8$. In these checks, as in GANDALF, $\Lambda$ and $\beta_i$ are independent inputs, and the identity holds for every pair. Not every pair corresponds to a plasma: by (\ref{eq:Lambdapm}), $\Lambda^\pm=\pm\sqrt5$ at $\beta_i=\tau=Z=1$, but no $\tau/Z>0$ gives $\Lambda=3$ at $\beta_i=0.37$. Through GANDALF's time step without hypercollisions, a state that starts in a single slice stays in it. With $M=64$, the copy closure, $\Lambda=-\sqrt5$, $\beta_i=2.3$, and GANDALF's damping of $z^\pm$ and $g$ at large $k_\perp$ and $k_z$ switched on, the amplitude that leaks into the other slices in 40 steps is at most $1.5\times10^{-14}$ of the initial one, for both time-stepping schemes and for initial slices at $u=0$ and at the largest phase velocity (figure~\ref{fig:decoupling}). In floating-point arithmetic the decoupling holds to round-off. Our checks are in double precision. GANDALF computes in single precision unless 64-bit arithmetic is enabled in JAX, and a run in single precision keeps the decoupling only to single-precision round-off.

The conservation laws of the slices need more than (\ref{eq:truncslice}). The step from (\ref{eq:truncslice}) to the conservation of $\int a_j^2\,\dd^3r$, or to the two-slice law (\ref{eq:twoslice}), uses the antisymmetry of the bracket under the volume average, $\int f\{h,f'\}\,\dd^3r=-\int f'\{h,f\}\,\dd^3r$. The dealiased pseudo-spectral bracket has this property exactly only when $\Phi$, $\Psi$ and $g$ lie inside the dealiasing band. GANDALF multiplies $g$ by the dealiasing mask after every step. It does not apply the mask to $z^\pm$, but every nonlinear term of their right-hand side is dealiased and their linear terms act mode by mode, so $z^\pm$ stay inside the band when their initial data and their forcing do.

Dissipation that multiplies every $g_m$ by the same factor, such as GANDALF's damping proportional to a power of $k_\perp^2$ or of $k_z^2$, commutes with the slice projection and damps each slice separately. The hypercollisions, which act on the Hermite index, couple the slices (section~\ref{tr:collsec}). Forcing of a moment loads the slices as in (\ref{eq:sliceforcing}). With $\lambda_j=1/n_j=R_{0j}^2$, an increment $f$ of $g_0$ adds $p_0\sqrt{\lambda_j}\,f$ to $a_j$, and an increment $f$ of $g_1$ adds $\sqrt2\,u_j\sqrt{\lambda_j}\,f$, because $r_0=1$ and $r_1(u)=\sqrt2\,u$. These are the continuum sources of section~\ref{sec:slices} with $\sqrt{\varrho(u)}$ replaced by $\sqrt{\lambda_j}$, which section~\ref{tr:gausssec} identifies, for the zero closure, as the square root of a quadrature weight. The projection acts mode by mode with real coefficients, so each slice receives, Fourier mode by Fourier mode, the pattern with which the moment is forced, including the modes on the plane $k_x=0$.

Parts of this structure are known. \citet[\S6.2.1]{Loureiro2016} diagonalised the truncated streaming matrix of the KREHM model to write the linear parallel step of Viriato in characteristic variables. There the matrix depends on $k_\perp$, and the characteristic fields are used only for that linear step. For the linearised Vlasov--Poisson system, \citet{ParkerDellar2015} identified the eigenvectors of the collisionless truncated Hermite system with discrete Case--Van Kampen modes and noted that the truncation is equivalent to a discretisation of velocity space on the roots of a Hermite polynomial. \citet{Tassi2015} constructed normal fields of truncated Hermite hierarchies obtained as fluid reductions of electromagnetic drift-kinetic equations. \citet[\S7]{Tassi2022} noted that the eigenvalues of the Hermite matrix are the zeros of a Hermite polynomial, a known property of this Jacobi matrix, and showed that the normal fields are proportional to the truncated distribution function at those velocities. In two dimensions each normal field is advected by a stream function of the form $\phi-\lambda_iA$, and in three dimensions the normal fields are coupled through source terms that contain the parallel gradients of the self-consistent potentials \citep[\S7]{Tassi2022}. The construction relies on the symmetry of the Hermite matrix \citep{Tassi2015}. In KRMHD the Landau coupling enters only the entry $A_{10}$ of a matrix that does not depend on the wavenumber and is symmetrised by $\Pw$, and the Alfv\'enic fields do not depend on $g$. The decoupling (\ref{eq:truncslice}) is therefore exact in three dimensions, for every $\Lambda\notin[0,1]$ and for both of GANDALF's closures.

\subsection{Gauss quadrature of the Case--Van Kampen measure}\label{tr:gausssec}

The polynomials $r_m$ connect the truncated hierarchy with the continuum. For the Hermite basis functions $e_m=H_mF_M/\sqrt{2^mm!}$ of section~\ref{sec:slices}, the slice transform is $\Tr[e_m](u)=\fe{F_u}{e_m}=p_m\,r_m(u)$.

\begin{proposition}\label{tr:gaussprop}
Let $\Lambda\notin[0,1]$ and use the zero closure. Then:
\begin{itemize}
\item[(i)] the polynomials $r_m$ are orthogonal for the spectral density, $\int\varrho\,r_m\,r_n\,\dd u=\delta_{mn}/p_m$;
\item[(ii)] the eigenvalues $u_j$ of $A$ are the zeros of $r_{M+1}$, and with $\lambda_j=1/n_j$
\begin{equation}\label{eq:gauss}
  \sum_{j=0}^M\lambda_j\,f(u_j)=\int\varrho(u)\,f(u)\,\dd u\qquad\text{for every polynomial $f$ of degree at most $2M+1$,}
\end{equation}
so that the $u_j$ and $\lambda_j$ are the nodes and weights (Christoffel numbers) of the $(M+1)$-point Gauss rule of the measure $\varrho\,\dd u$, and $\sum_j\lambda_j=1/p_0$;
\item[(iii)] $\boldsymbol r_j$ holds the first $M+1$ Hermite coefficients of $F_{u_j}$, and for $G=\sum_{m\le M}g_me_m$ the discrete slice amplitude is the continuum transform at the node, $\boldsymbol w_j^{\mathsf T}g=\Tr[G](u_j)$, so that $a_j=\sqrt{\lambda_j}\,\Tr[G](u_j)$.
\end{itemize}
\end{proposition}

\begin{proof}
(i) By theorem~\ref{thm:spectrum} the map $G\mapsto a$ is unitary, and the slice amplitude of $e_m$ is $\sqrt\varrho\,p_mr_m$. Hence $p_m\delta_{mn}=\fe{e_m}{e_n}=\int\varrho\,p_mr_m\,p_nr_n\,\dd u$. (ii) The polynomials $\sqrt{p_m}\,r_m$ are therefore orthonormal for $\varrho\,\dd u$. Their three-term recurrence, which follows from $LF_u=uF_u$ (appendix~\ref{app:truncated}), has the symmetric Jacobi matrix $\Pw^{1/2}A\,\Pw^{-1/2}$, and truncation with the zero closure keeps its first $M+1$ rows and columns. By the theorem of \citet{GolubWelsch1969}, the eigenvalues of this truncated matrix are the zeros of $r_{M+1}$ and the nodes of the Gauss rule, and the weights are $\mu_0O_{0j}^2$, where $O$ holds its orthonormal eigenvectors and $\mu_0=\int\varrho\,\dd u=1/p_0$. Since $O_{mj}=\sqrt{p_m}\,R_{mj}$, the weights are $R_{0j}^2=1/n_j$. Exactness for degree $2M+1$ is Gauss's theorem. (iii) The first statement is the definition of $r_m$. For $G=\sum_mg_me_m$, $\Tr[G](u_j)=\sum_mg_m\,p_mr_m(u_j)=\boldsymbol w_j^{\mathsf T}g$.
\end{proof}

The continuum spectral density $\varrho\,\dd u$ is thus the orthogonality measure of the ladder matrix, and truncation with the zero closure replaces the continuum of slices by the Gauss rule of that measure. The weights are close to the measure of the interval around each node, $\lambda_j\simeq\varrho(u_j)\Delta_j$, with $\Delta_j$ the local spacing of the nodes; on $|u_j|<3$ the relative difference is at most $4.5\times10^{-4}$ for $\Lambda=\sqrt5$ and $5.4\times10^{-5}$ for $\Lambda=-\sqrt5$ at $M=1024$. By proposition~\ref{tr:gaussprop}(iii), $a_j\simeq\Delta_j^{1/2}\,a_{u_j}$ for a state in the span of the retained moments, with $a_u$ the continuum slice amplitude (\ref{eq:aunit}): the discrete slices are the continuum slices at the nodes, weighted by the square root of the node spacing. The conserved quadratic forms of section~\ref{tr:formsec}, $\sum_j\varphi(u_j)\,a_j^2$, are the Gauss approximations of $\fe{G}{\varphi(L)G}=\int\varphi(u)\,a_u^2\,\dd u$. For a smooth state outside the span of the retained moments (in our test a sum of two shifted Gaussians), the distribution of free energy over the slices converges to the continuum one. The largest difference between the cumulative distributions falls from $1.4\times10^{-2}$ of the total free energy at $M=16$ to $2.1\times10^{-4}$ at $M=1024$ for $\Lambda=\sqrt5$, approximately as $M^{-1}$, and from $8.8\times10^{-3}$ to $1.8\times10^{-4}$ for $\Lambda=-\sqrt5$. Figure~\ref{fig:spectrum} shows $\varrho$ with the nodes and weights.

Landau damping in a truncated hierarchy is a Gauss quadrature of the continuum phase mixing. For $G(0)=F_M$ and a single Fourier mode in a uniform field, $a_j(0)=p_0\sqrt{\lambda_j}$ and each slice rotates at its own frequency, so the density moment is $g_0(\hat t)=p_0\sum_j\lambda_je^{-\ii\hat t u_j}$ with $\hat t=sk_zt$. In the continuum the same argument gives $M_0[G](\hat t)=p_0\int\varrho(u)\,e^{-\ii\hat t u}\,\dd u$, whose decay at late times is a sum over the Landau roots. The two agree to $3\times10^{-15}$ for $\hat t\le12$ at $M=128$ and at $M=512$, for both signs of $\Lambda$, well before the recurrence produced by the discreteness of the nodes.

At $\Lambda=\infty$ proposition~\ref{tr:gaussprop} reduces to the known correspondence between the truncated Hermite hierarchy and Gauss--Hermite quadrature: the eigenvalues are the zeros of $H_{M+1}$, as \citet{ParkerDellar2015} and \citet{Tassi2022} noted, and the identification of the nodes and weights with the eigenvalues and eigenvectors of the Jacobi matrix is that of \citet{GolubWelsch1969}. For finite $\Lambda$ the measure is the dressed spectral density $\varrho=F_M/|\eps|^2$ rather than $F_M$, the nodes are the zeros of $r_{M+1}$ rather than of $H_{M+1}$, and the eigenvectors are the Hermite coefficients of the Case--Van Kampen modes. The identification of the truncated KRMHD hierarchy with the Gauss rule of the dressed measure at finite $\Lambda$, together with the exact nonlinear decoupling of theorem~\ref{thm:trunc}, is new; for linearised Vlasov--Poisson, \citet{ParkerDellar2015} had identified the eigenvectors with discrete Case--Van Kampen modes. The copy closure does not give the Gauss rule of $\varrho\,\dd u$: its eigenvalues are the zeros of $r_{M+1}-r_{M-1}$ (appendix~\ref{app:truncated}).

\subsection{Conserved quadratic forms}\label{tr:formsec}

The truncated counterpart of the family (\ref{eq:Qphi}) is
\begin{equation*}
  Q^{(M)}_\varphi=\int\dd^3r\;g^{\mathsf T}\Pw\,\varphi(A)\,g=\sum_{j=0}^M\varphi(u_j)\int\dd^3r\,a_j^2 ,
\end{equation*}
where the second form follows from $\Pw\varphi(A)=\Pw R\,\varphi(\mathrm{diag}\,u)\,R^{\mathsf T}\Pw$. The following proposition states that these are all the conserved quadratic forms whose matrix does not depend on position or wavenumber.

\begin{proposition}\label{tr:formsprop}
Let $\Lambda\notin[0,1]$, use either closure, and let $\Phi$, $\Psi$ and $g$ lie inside the dealiasing band, with no collisions, dissipation or forcing acting on $g$. A form $\int\dd^3r\,g^{\mathsf T}Kg$ with a real symmetric $(M+1)\times(M+1)$ matrix $K$ is conserved by (\ref{tr:hierarchy}) for all $\Phi$ and $\Psi$ if and only if $KA=A^{\mathsf T}K$, that is, if and only if $K=\Pw\varphi(A)$ for a real function $\varphi$ on the nodes. The conserved forms make up the $(M+1)$-dimensional family $Q^{(M)}_\varphi$, spanned by the matrices $\Pw A^n$ with $n=0,\dots,M$.
\end{proposition}

\begin{proof}
For band-limited fields the bracket and $\gradpar$ are antisymmetric under the volume average, $\int f\gradpar h\,\dd^3r=-\int h\gradpar f\,\dd^3r$. Hence $\int g^{\mathsf T}K\{\Phi,g\}\,\dd^3r=0$ for every symmetric $K$, and the $\boldsymbol E\times\boldsymbol B$ term conserves every such form. The streaming term gives $\dd/\dd t\int g^{\mathsf T}Kg\,\dd^3r=-2s\int g^{\mathsf T}KA\,\gradpar g\,\dd^3r$. By the antisymmetry of $\gradpar$, the symmetric part of $KA$ does not contribute, and only the antisymmetric part $B=\frac12(KA-A^{\mathsf T}K)$ remains. If $B_{mn}\neq0$ for some $m\neq n$, the state with $\Psi=0$, $g_m=\cos k_zz$, $g_n=\sin k_zz$ and all other moments zero gives $\int g^{\mathsf T}B\,\partial_zg\,\dd^3r=k_zB_{mn}V$, with $V$ the volume, so the form is not conserved. Hence $KA=A^{\mathsf T}K$. In the basis of eigenvectors this says that $R^{\mathsf T}KR$ commutes with $\mathrm{diag}(u_j)$, which has distinct entries, so $R^{\mathsf T}KR=\mathrm{diag}(\varphi_j)$ and $K=\Pw R\,\mathrm{diag}(\varphi_j)R^{\mathsf T}\Pw=\Pw\varphi(A)$ with $\varphi(u_j)=\varphi_j$. A function on $M+1$ distinct nodes coincides there with a polynomial of degree at most $M$, which gives the spanning set. Conversely, each $\int a_j^2\,\dd^3r$ is conserved by (\ref{eq:truncslice}) and the antisymmetry of the bracket and of $\partial_z$.
\end{proof}

All members of the family are conserved on GANDALF's right-hand side, at relative rates of at most $3.2\times10^{-17}$, for $M=9$ and $10$, $\Lambda=\pm\sqrt5$, both closures and brackets that dominate the streaming (figure~\ref{fig:invariants}). The restriction to matrices that do not depend on the wavenumber is necessary on the modes with $k_\perp\neq0$; the modes with $k_\perp=0$, which the brackets do not couple to the others, may carry a different member of the family. For a matrix $K_{\boldsymbol k}$ that depends on $\boldsymbol k$, let $\hat K$ be the operator with this symbol. The $\boldsymbol E\times\boldsymbol B$ term then changes the form at the rate $-2\int\Phi\sum_m\{g_m,(\hat Kg)_m\}\,\dd^3r$, which vanishes for every $\Phi$ only if $\sum_m\{g_m,(\hat Kg)_m\}=0$ for every $g$. For a single non-zero moment $g_m=\cos(\boldsymbol k\cdot\boldsymbol r)+\cos(\boldsymbol q\cdot\boldsymbol r)$ this sum is $(c_{\boldsymbol q}-c_{\boldsymbol k})\{\cos\boldsymbol k\cdot\boldsymbol r,\cos\boldsymbol q\cdot\boldsymbol r\}$, where $c_{\boldsymbol k}$ is the symbol of $\hat K_{mm}$, so the symbol must be the same at any two wavevectors with non-parallel perpendicular parts; this is the triad argument of section~\ref{sec:invariants}. On a test state the free energy weighted by $|k_z|$ or by $k_\perp^2$ changes at relative rates of $2.5\times10^{-3}$ and $4.5\times10^{-3}$, against $2\times10^{-18}$ for the free energy itself.

The members $\varphi=1$ and $\varphi=u$ are CMM's invariants on the retained moments. For the zero closure and $G=\sum_{m\le M}g_me_m$, $\fe{G}{G}=g^{\mathsf T}\Pw g$ and $\fe{G}{LG}=g^{\mathsf T}\Pw Ag$, so that, by (\ref{eq:WGH}), $W^\pm=(n_iT_i/2V)\,Q^{(M)}_1$ and $\Gamma^\pm=(\vth/2V)\,Q^{(M)}_u$ are CMM's (B31) and (B32) for the truncated state. For $\varphi=u^2$ and higher powers the truncated form misses a boundary term at $m=M$; for example, $\fe{G}{L^2G}=g^{\mathsf T}\Pw A^2g+\frac12(M+1)\,g_M^2$. Under the copy closure, $Q^{(M)}_1$ carries the weight (\ref{tr:pM}), about one half, on $g_M^2$, and it is not (B31). GANDALF's diagnostic \texttt{hermite\_free\_energy} uses the weights $p_0=1-1/\Lambda$ and $1$ for every $m\ge1$, including $m=M$. It is conserved under the zero closure but not under the copy closure, for which we measured relative rates of change between $4\times10^{-4}$ and $3\times10^{-3}$ on GANDALF's right-hand side, against $3\times10^{-17}$ for $Q^{(M)}_1$. The form $Q^{(M)}_u$ is the same for both closures, because $\Pw A$ is. The non-quadratic slice invariants (\ref{eq:casimir}) do not survive the Fourier--Galerkin truncation in $\boldsymbol r$: evaluated with exact quadrature on a grid padded by a factor of two, $\int a_j^2\,\dd^3r$ is conserved on GANDALF's right-hand side to $9\times10^{-17}$, while $\int a_j^4\,\dd^3r$ changes at relative rates of $4.9\times10^{-4}$ and more.

\subsection{The phase-space helicity at odd and even \texorpdfstring{$M$}{M}}\label{tr:helsec}

CMM's helicity (\ref{eq:WGH}) has the kernel $\PVint(\cdot)^2/(uF_M)\,\dd u$, which is not banded in the Hermite basis. Its matrix on the retained moments, $(B_M)_{mn}=\PVint\psi_m\psi_n\,u^{-1}\,\dd u$ for $m,n\le M$, has to be compared with the conserved forms $\Pw\varphi(A)$ with $\varphi(u)=1/u$. Let $J$ be the ladder matrix of the zero closure at $\Lambda=\infty$, which is the Jacobi matrix of the polynomials $h_m(u)=H_m(u)/\sqrt{2^mm!}$ (these are the $r_m$ at $\Lambda=\infty$), let $\lambda^\infty_j$ be the weights of the $(M+1)$-point Gauss--Hermite rule, normalised so that they add up to one, and let $\boldsymbol h(0)$ and $\boldsymbol h'(0)$ be the vectors with components $h_m(0)$ and $h_m'(0)$. For the zero closure,
\begin{equation}\label{eq:Hoddeven}
  B_M=J^{-1}=\Pw A^{-1}\quad(M\ \text{odd}),\qquad
  B_M=J^{+}+\lambda^\infty_{M/2}\big[\boldsymbol h(0)\,\boldsymbol h'(0)^{\mathsf T}+\boldsymbol h'(0)\,\boldsymbol h(0)^{\mathsf T}\big]\quad(M\ \text{even}),
\end{equation}
where $J^+$ is the Moore--Penrose inverse of $J$ and $\lambda^\infty_{M/2}$ is the weight of the node $u=0$. Appendix~\ref{app:truncated} proves (\ref{eq:Hoddeven}) from the exactness of Gauss--Hermite quadrature.

For odd $M$, $A$ is invertible for every $\Lambda\neq1$, and $\Pw A^{-1}=J^{-1}$ because $A=J\Pw$ for the zero closure. The conserved form $Q^{(M)}_{1/u}=\sum_ju_j^{-1}\int a_j^2\,\dd^3r$ is then exactly CMM's helicity of the truncated state, $\Hph^\pm=(n_iT_i/2\vth V)\,Q^{(M)}_{1/u}$, and its matrix does not depend on $\Lambda$. This holds for the zero closure. Under the copy closure the conserved form at odd $M$ is $\int g^{\mathsf T}D_cJ^{-1}D_cg\,\dd^3r$, whose matrix differs from $B_M$ by 33\% at $M=9$ and 23\% at $M=33$, relative to the largest element.

For even $M$, $u=0$ is an exact eigenvalue, $u_{M/2}=0$, for every $\Lambda$ and both closures. The reason is the parity of the ladder: $SAS=-A$ with $S=\mathrm{diag}((-1)^m)$, so the spectrum is symmetric about zero, and a matrix of odd dimension with a symmetric spectrum has the eigenvalue zero. For the zero closure the left null vector is $p_0\boldsymbol h(0)$ for every $\Lambda$, so the slice at $u=0$ is proportional to $\sum_mh_m(0)\,g_m=G_M(0)/F_M(0)$, the truncated distribution $G_M=\sum_{m\le M}g_me_m$ at $v_\parallel=0$; for the copy closure the last component of the null vector is multiplied by $p_M$. By (\ref{eq:truncslice}) this slice is advected by $\uE$ alone, $\partial_ta_{M/2}+\{\Phi,a_{M/2}\}=0$, for both closures, and its energy $\int a_{M/2}^2\,\dd^3r$ is conserved; we checked both properties on GANDALF's right-hand side, to $2\times10^{-15}$ and $2\times10^{-16}$ (appendix~\ref{app:truncated}). For the zero closure its energy is proportional to the imaginary part of the complex helicity (\ref{eq:Hcomplex}) of the truncated state, so that its conservation is the truncated form of the conservation of CMM's $\mathrm{Im}\,\Hph$; for the copy closure the conserved slice energy differs from it through the weight $p_M$ on $g_M$. CMM showed that in full gyrokinetics the function $\psi_s$ of their eq.~(2.18), evaluated at $v_\parallel=0$, is advected by the gyroaveraged generalised potential $\langle\chi\rangle_{\boldsymbol R_s}$ alone (their eq.~2.22), which gives the conservation of $\mathrm{Im}\,\Hph$ (their eq.~2.23). At odd $M$ there is no such slice, because $\det A\neq0$ (appendix~\ref{app:truncated}); the vector $\boldsymbol h(0)$ then fails to be a null vector only through the component $h_{M-1}(0)\sqrt{M/2}$ of $\boldsymbol h(0)^{\mathsf T}A$, in column $M$.

At $\Lambda=\infty$ the even-$M$ matrix $B_M$ in (\ref{eq:Hoddeven}) is not of the form $\Pw\varphi(A)$, because its rank-two term couples the null vector $\boldsymbol h(0)$ of $J$ to its other eigenvectors, so by proposition~\ref{tr:formsprop} the restriction of $\Hph$ to the retained moments is not conserved. At finite $\Lambda$ we found the same numerically: on GANDALF's right-hand side at $M=10$ the restriction changes at relative rates of $2\times10^{-4}$ and $2\times10^{-3}$ for $\Lambda=\sqrt5$ and $-\sqrt5$. The conserved principal-value form is
\begin{equation*}
  Q^{(M)}_{1/u}=\sum_{u_j\neq0}\frac{1}{u_j}\int a_j^2\,\dd^3r=\int\dd^3r\;g^{\mathsf T}\Pw A^{D}g,
\end{equation*}
with $A^D$ the Drazin inverse. It omits the slice at $u=0$. At $\Lambda=\infty$ its integrand differs from $g^{\mathsf T}B_Mg$ by exactly $-2\lambda^\infty_{M/2}\,[\boldsymbol h(0)\cdot g]\,[\boldsymbol h'(0)\cdot g]$, which is the node term that a principal-value integral drops from the Gauss rule. Since $\lambda^\infty_{M/2}\simeq[\pi/(2M+3)]^{1/2}$, the even-$M$ form converges to the helicity only as $M^{-1/2}$.

The difference is large for a smooth state. For the shifted Maxwellian $G=F_M(u-\frac12)$, all of whose Hermite coefficients are non-zero, $Q^{(M)}_1$ and $Q^{(M)}_u$ equal their continuum values to $4\times10^{-16}$ from $M=16$ on, and the odd-$M$ helicity has errors of $4.5\times10^{-6}$, $4.5\times10^{-12}$ and $2\times10^{-16}$ at $M=9$, 17 and 33, for each $\Lambda$ tested. The error of the even-$M$ principal-value form times $M^{1/2}$ tends to $-2.86$, $-1.01$ and $-1.52$ for $\Lambda=\sqrt5$, $-\sqrt5$ and $\infty$, with a relative correction of order $M^{-1/2}$ at finite $\Lambda$ and of order $M^{-1}$ at $\Lambda=\infty$. At $M=64$ and 128 the error is between 5\% and 18\% of the helicity, whose continuum value for this state is $\PVint G^2(uF_M)^{-1}\,\dd u=2e^{1/2}D(1)=1.774$, with $D$ the Dawson function. Diagnostics of the helicity should therefore use odd $M$ and the zero closure, for which the conserved form is CMM's helicity on the retained moments. At even $M$ the principal-value form is conserved but carries an error of order $M^{-1/2}$, while the slice at $u=0$, and with it $\mathrm{Im}\,\Hph$, is exact.

The field advected by the $\boldsymbol E\times\boldsymbol B$ flow alone at an odd number of retained moments, and its family of Casimir invariants in three dimensions, were found by \citet{Tassi2022} for Hermite truncations of electromagnetic drift-kinetic models with a symmetric matrix \citep[see also][]{Tassi2023}; the two-dimensional version is in \citet{Tassi2015}. What is added here is that this field survives the $\Lambda$-modified ladder and both of GANDALF's closures, with a null vector that does not depend on $\Lambda$, its identification, for the zero closure, with CMM's $\mathrm{Im}\,\Hph$, and the consequence (\ref{eq:Hoddeven}) for the real part of the helicity.

\subsection{The phase velocities}\label{tr:velsec}

At $\Lambda=\infty$ and with the zero closure, $A=J$, its eigenvalues are the Gauss--Hermite nodes $x_j$, the zeros of $H_{M+1}$, and the $\lambda_j$ are the Gauss--Hermite weights $\lambda^\infty_j$. A finite $\Lambda$ changes only the entry $A_{10}$, and expansion of the determinant along the first row gives the exact identity
\begin{equation}\label{tr:det}
  \det(z-A)=\det(z-J)+\frac{1}{2\Lambda}\det\big(z-J^{(2)}\big),
\end{equation}
where $J^{(2)}$ is $J$ without its first two rows and columns; the same identity holds for the copy closure, with $J$ replaced by the copy-closure matrix at $\Lambda=\infty$. The traces of $A$ and $A^2$ give the sum rules $\sum_ju_j=0$ and $\sum_ju_j^2=\frac12M(M+1)-1/\Lambda$, with $\sqrt{M(M+1)}$ added for the copy closure. To first order in $1/\Lambda$, $u_j=x_j-\lambda^\infty_jx_j/\Lambda$ (appendix~\ref{app:truncated}). At $\Lambda=\pm\sqrt5$ the first-order formula is only qualitative: at $M=128$ the largest shift is $0.125$ of the local spacing for $\Lambda=\sqrt5$ and $0.097$ for $\Lambda=-\sqrt5$, and the formula misses it by $0.032$ and $0.020$ of the spacing. For large $M$ the shifts follow the spectral-shift law
\begin{equation}\label{tr:shift}
  \frac{u_j-x_j}{\Delta_j}\to\frac{1}{\pi}\arg\eps(x_j+\ii0)\qquad(M\to\infty),
\end{equation}
with $\Delta_j$ the local spacing of the nodes and $\eps$ the dielectric function (\ref{eq:eps}). We have not proved (\ref{tr:shift}). We state it as an asymptotic law, which we verified numerically: on $|x_j|<3$ the largest deviation from (\ref{tr:shift}) is $4.7\times10^{-3}$, $2.4\times10^{-3}$, $1.2\times10^{-3}$ and $6.1\times10^{-4}$ at $M=128$, 512, 2048 and 8192 for $\Lambda=\sqrt5$, and $1.7\times10^{-3}$, $8.5\times10^{-4}$, $4.3\times10^{-4}$ and $2.1\times10^{-4}$ for $\Lambda=-\sqrt5$, decreasing approximately as $M^{-1/2}$. Appendix~\ref{app:truncated} gives the argument that leads to (\ref{tr:shift}). Since $\mathrm{Im}\,\eps(u+\ii0)=-\pi uF_M(u)/\Lambda$, the branch $\Lambda>0$ draws the slices toward $u=0$ and the branch $\Lambda<0$ pushes them away: at $M=128$ the spacing at $u=0$ is $0.1807$ for $\Lambda=\sqrt5$, $0.1952$ for $\Lambda=\infty$ and $0.2020$ for $\Lambda=-\sqrt5$.

The largest phase velocity $u_{\max}=u_M$ hardly depends on $\Lambda$. The normalised extreme eigenvector of $J$ has the components $(\lambda^\infty_M)^{1/2}h_m(x_M)$, so its components on the moments $m=0$ and $1$, where $\Lambda$ acts, are of order $(\lambda^\infty_M)^{1/2}$ up to an algebraic factor. The first-order shift $-\lambda^\infty_Mx_M/\Lambda$ is of size $x_Me^{-x_M^2}/\Lambda$ up to an algebraic factor, with $x_M^2=2M+3-2^{2/3}|\alpha_1|(2M+3)^{1/3}+\dots$ and $\alpha_1$ the first zero of the Airy function; at $M=64$ it is about $10^{-49}$. In double precision $u_{\max}$ is the same for $\Lambda=\sqrt5$, $-\sqrt5$ and $\infty$ to $3\times10^{-13}$ for $M\ge32$. Up to such terms, $u_{\max}$ is the largest zero of the Hermite function $\psi_{M+1}$ for the zero closure, and the largest root of $\psi_{M+1}=\psi_{M-1}$ for the copy closure. With $\kappa_m=(2m+1)^{1/2}$ as in section~\ref{sec:decorrelation}, appendix~\ref{app:truncated} derives
\begin{equation}\label{eq:umax}
\begin{aligned}
  u_{\max}&=\kappa+2^{-1/3}\alpha_1\,\kappa^{-1/3}-\frac{2^{-2/3}\alpha_1^2}{10}\,\kappa^{-5/3}+O(\kappa^{-3}),&\qquad\kappa&=\kappa_{M+1}\quad\text{(zero closure)},\\
  u_{\max}&=\kappa+2^{-1/3}\alpha_1'\,\kappa^{-1/3}-2^{-5/3}\Big(\frac{4}{5\alpha_1'}+\frac{\alpha_1'^2}{5}\Big)\kappa^{-5/3}+O(\kappa^{-3}),&\qquad\kappa&=\kappa_M\quad\text{(copy closure)},
\end{aligned}
\end{equation}
where $\alpha_1=-2.3381074$ and $\alpha_1'=-1.0187930$ are the first zeros of the Airy function and of its derivative. For the zero closure the first two terms are the classical asymptotics of the largest zero of a Hermite polynomial, and the third follows from the same uniform Airy approximation; the copy-closure line follows from the Robin condition derived in appendix~\ref{app:truncated}. The coefficients of $\kappa^{-1/3}$ and $\kappa^{-5/3}$ are $-1.855757$ and $-0.344383$ for the zero closure and $-0.808617$ and $+0.181950$ for the copy closure. Values of $u_{\max}$ computed up to $M=16384$ confirm these coefficients, and fits to them give about $-0.169$ and $-0.144$ for the coefficients of $\kappa^{-3}$, which we have not derived. At $M=64$ and 128 the zero-closure formula gives $10.616139$ and $15.355091$, against the exact $10.616023$ and $15.355050$. Table~\ref{tr:umaxtab} lists $u_{\max}$ for GANDALF's matrix, and figure~\ref{fig:spectrum} shows it as a function of $M$. The copy closure has the larger $u_{\max}$.

\begin{table}[t]
  \centering
  \caption{The largest phase velocity $u_{\max}$ of the truncated hierarchy, the largest eigenvalue of GANDALF's ladder matrix. The values are the same for $\Lambda=\sqrt5$, $-\sqrt5$ and $\infty$ to all digits shown.}
  \label{tr:umaxtab}
  \begin{tabular}{@{}rcc@{}}
    \toprule
    $M$ & zero closure & copy closure \\
    \midrule
     16 &  4.87135 &  5.30215 \\
     32 &  7.25385 &  7.66432 \\
     64 & 10.61602 & 11.00116 \\
     65 & 10.70522 & 11.08975 \\
    128 & 15.35505 & 15.71228 \\
    129 & 15.41803 & 15.77494 \\
    256 & 22.03623 & 22.36470 \\
    \bottomrule
  \end{tabular}
\end{table}

\subsection{Hypercollisions}\label{tr:collsec}

GANDALF damps the moments with $m\ge2$ by hypercollisions, $\partial_tg_m=-\nu(m/M)^ng_m$. In slice variables this is
\begin{equation}\label{tr:coll}
  \partial_ta_l\big|_{\mathrm{coll}}=-\sum_j\mathsf C_{lj}\,a_j,\qquad
  \mathsf C_{lj}=\nu\sum_{m=2}^Mp_m\Big(\frac mM\Big)^nR_{ml}\,R_{mj},
\end{equation}
which is $\nu R^{\mathsf T}\Pw\mathcal DR$ with $\mathcal D=\mathrm{diag}((m/M)^n)$ for $m\ge2$ and $\mathcal D_{00}=\mathcal D_{11}=0$, the truncated form of the collision kernel of (\ref{eq:sliceforcing}). The matrix $\mathsf C$ is symmetric and positive semi-definite. It is the only term of GANDALF's time step that couples different slices, and its off-diagonal elements are not small compared with the diagonal ones (section~\ref{sec:slices}). Through GANDALF's time step, a state that starts in one slice leaks into the others at first order in $\nu$: doubling $\nu$ doubles the leaked amplitude (the measured ratio is $1.9999$), and the leaked amplitude agrees with the first-order prediction from (\ref{tr:coll}) to better than 1\%, for both schemes and both closures.

The diagonal elements $\mathsf C_{jj}$ are the damping rates of the slices. For $|u_j|\ll(2M)^{1/2}$ the eigenvector spreads its weight over $0\le m\le M$ as $R_{mj}^2\simeq1/[2(Mm)^{1/2}]$, on average over neighbouring $m$, and
\begin{equation}\label{tr:gamma}
  \mathsf C_{jj}\simeq\frac{\nu}{2M^{1/2}}\int_0^M\Big(\frac mM\Big)^n\frac{\dd m}{m^{1/2}}=\frac{\nu}{2n+1}.
\end{equation}
At $\Lambda=\sqrt5$ the slice at $u=0$ has $\mathsf C_{jj}=0.074\,\nu$ for $n=6$ ($M=64$ and $128$), against $\nu/13=0.077\,\nu$, and $0.18\,\nu$ to $0.19\,\nu$ for $n=2$, against $0.2\,\nu$. The fastest slices are damped much more strongly, at $0.63\,\nu$ to $0.73\,\nu$ for $n=6$ and $0.85\,\nu$ to $0.90\,\nu$ for $n=2$ ($M=64$ and $128$), because their eigenvectors are concentrated at large $m$. At $\Lambda=\infty$ the ratio $(2n+1)\mathsf C_{jj}/\nu$ for the slices nearest $u=0$ approaches 1 as $M^{-1}$. At $\Lambda=\sqrt5$ the ratio is $1.049$ at $M=64$ and $0.981$, $0.978$ and $0.982$ at $M=256$, 1024 and 2048 ($n=6$); we have not derived its behaviour at larger $M$.

\subsection{Reflection at the closure}\label{tr:reflsec}

Both closures are conservative, because $\Pw A$ is symmetric for both: no free energy leaves the hierarchy at $m=M$, and a phase-mixing packet that reaches $m=M$ is reflected. Consider the linear problem at a single wavenumber $k$ along the field, with hypercollisions of order $n$. By theorem~\ref{thm:frozen} this problem also describes phase mixing along field lines frozen by ideal induction. A packet launched from $g_0$ returns to $g_0$, and the ratio $\mathcal R$ of its returning amplitude to the amplitude in the same system without hypercollisions is, for large $M$, given approximately by
\begin{equation}\label{eq:reflection}
  \mathcal R=e^{-X},\qquad X=\frac{2\nu t_w}{2n+1},\qquad t_w=\frac{\sqrt{2M}}{|k|s} .
\end{equation}
In the ray derivation, free streaming from $G(0)=F_M$ at $\Lambda=\infty$ makes $|g_m|^2$ a Poisson distribution in $m$ with mean $\bar m=(kst)^2/2$, so the packet reaches $m=M$ at $t=t_w$. On the way it is damped at the local rate $\nu(\bar m/M)^n=\nu(t/t_w)^{2n}$, whose integral over $0<t<t_w$ is $\nu t_w/(2n+1)$, and the return costs the same. In the mode derivation, the revival is the re-phasing of the slices near $u=0$, which happens after a time close to $2t_w$, and these slices are damped at the rate (\ref{tr:gamma}). Appendix~\ref{app:truncated} gives both derivations.

We compared (\ref{eq:reflection}) with the exact revival computed from GANDALF's linear operator in 360 cases: $M$ from 32 to 256, $n=2$, 4 and 6, $\nu$ from $0.3$ to 30, three wavenumbers, $\Lambda=\infty$ and $\sqrt5$, and $X$ between $0.015$ and 43. In all 319 cases in which the revival stands out from the non-recurrent response, the law overestimates $\mathcal R$. The relative error of the exponent is a finite-$M$ effect, with a median of $0.139$, $0.068$, $0.034$ and $0.017$ at $M=32$, 64, 128 and 256, and for $M\ge128$ it is at most $5.0\%$ at every $X$ tested. GANDALF's own time step reproduces the exact revival: for $M=128$, $n=6$, $\nu=3$, $\Lambda=\sqrt5$ and $k_zs=2\pi$ the ratio is $0.2913$, against $e^{-X}=0.3087$ (figure~\ref{fig:reflection}). The revival occurs near $2t_w[(2M+3)/2M]^{1/2}$ at every $\Lambda$, with a shift of order $M^{-1/2}$ at finite $\Lambda$ (appendix~\ref{app:truncated}). The ratio itself hardly depends on $\Lambda$: in the case above it is $0.2913$ at $\Lambda=\sqrt5$ and $0.2910$ at $\Lambda=\infty$. This is consistent with the ray picture, in which the attenuation is accumulated at $m\ge2$, where the hypercollisions act and the rows of the ladder matrix do not depend on $\Lambda$. The coupling of $g_0$ to $g_1$, the only place where $\Lambda$ enters, affects the launch and the return of the packet in the same way with and without hypercollisions.

Reflection changes forced steady states. With white-noise forcing of $g_0$, every packet bounces between the low moments and the closure, and it loses the factor $e^{-2X}$ in energy per round trip. Summing the forward and backward passes, with $\mu=(m/M)^{n+1/2}$ (appendix~\ref{app:truncated}),
\begin{equation}\label{tr:coth}
  \frac{W(m)}{W_{\mathrm{open}}(m)}=\frac{1+e^{-2X(1-\mu)}}{1-e^{-2X}}\to\coth X\qquad(m\ll M),
\end{equation}
where $W_{\mathrm{open}}(m)$ is the spectrum of the same system with nothing returning from $m=M$. For $M=128$, $n=6$, $\nu=3$, $\Lambda=\sqrt5$ and $k_zs=2\pi$, the law (\ref{eq:reflection}) gives $X=1.175$ and $\coth X=1.211$, and the exact revival gives $X=1.234$ and $\coth X=1.185$. Steady states computed with GANDALF's linear operator (Lyapunov solves) give $W/W_{\mathrm{open}}=1.1852$ to $1.1911$ for $2\le m\le48$, whether the open reference is built with an outgoing closure or with an extended hierarchy, and (\ref{tr:coth}) with the measured $X$ matches them to $0.5\%$. A reflecting closure at these parameters therefore raises the Hermite spectrum at low $m$ by 18.5\% above the spectrum without reflection. At larger $m$ the law underestimates the effect of the closure: the computed ratio is $1.206$ to $1.207$ at $m=64$ and $1.400$ to $1.407$ at $m=96$, for the two open references, against $1.188$ and $1.228$ from (\ref{tr:coth}). The reflection law and its steady-state form are new. They were tested only for linear problems at a single wavenumber; we have not tested them along perturbed field lines or with nonlinear transfer in $k_\perp$.

\subsection{Consequences for time stepping and diagnostics}\label{tr:numsec}

The slice decomposition also describes what an explicit time step does to the hierarchy. Both of GANDALF's schemes integrate the brackets explicitly. The explicit term is $N(g)=-\{\Phi,g\}-s\{\Psi,Ag\}$, and by theorem~\ref{thm:trunc} its projection on slice $j$ is $-\{\Phi+su_j\Psi,a_j\}$. Slice $j$ is therefore advected explicitly at the speed
\begin{equation*}
  V_j=\max_{\boldsymbol r}\big|\nabla_\perp(\Phi+su_j\Psi)\big|\le\max|\nabla_\perp\Phi|+s|u_j|\max|\nabla_\perp\Psi| .
\end{equation*}
When $s|u_j|\max|\nabla_\perp\Psi|\gg\max|\nabla_\perp\Phi|$ the last term dominates, and the fastest slice moves at up to about $su_{\max}\max|\nabla_\perp\Psi|$, with $u_{\max}=10.6$ at $M=64$ and $15.4$ at $M=128$ (table~\ref{tr:umaxtab}). GANDALF's time-step estimate uses the speed $\max(v_A,|\nabla_\perp z^\pm|,s|\nabla_\perp\Psi|)$, which omits the factor $u_{\max}$. An earlier GPU version of GANDALF reduced its time step by the factor $M^{-1/2}$ \citep[eq.~2.26]{KanekarThesis2015}, which has the scaling of $u_{\max}\simeq(2M)^{1/2}$; version 0.6.0 does not. \citet{Mandell2024} include the largest Gauss--Hermite node in the parallel-streaming frequency that sets the time step of the GX code; the perpendicular advection of the fastest slice by the field-line term, at the speed $su_{\max}|\nabla_\perp\Psi|$, does not enter their estimate.

The quantity that limits the step is the spectral radius $\rho_j$ of the discrete advection operator $a\mapsto-\{\Phi+su_j\Psi,a\}$ on the dealiasing band, which satisfies $\rho_j\le k_{\perp,\max}V_j$, with $k_{\perp,\max}$ the largest perpendicular wavenumber of the band. For band-limited fields this operator is antisymmetric, so its eigenvalues are imaginary. The integrating-factor scheme Lawson-RK4 \citep{Lawson1967} applies classical RK4 to the explicit term. Its stability function satisfies $|R_4(\ii y)|^2=1-y^6/72+y^8/576$, so it has a true limit: the explicit part is stable when $\Delta t\,\rho_j\le2\sqrt2$ for every slice, a condition set by the fastest one. The explicit stages of the IMEX-RK222 scheme, which is ARS(2,2,2) \citep{AscherRuuthSpiteri1997}, have the stability function $R_2(\zeta)=1+\zeta+\zeta^2/2$ and
\begin{equation*}
  |R_2(\ii y)|^2=1+\frac{y^4}{4}>1\qquad(y\neq0),
\end{equation*}
so they have no stable interval on the imaginary axis, and their stability rests on dissipation. If the modes of slice $j$ that rotate fastest are damped at the rate $d_j$, the condition is approximately
\begin{equation}\label{eq:explicitstability}
  \frac{(\Delta t\,\rho_j)^4}{8}\le\Delta t\,d_j ,
\end{equation}
where $d_j$ is the hyper-resistive damping rate of those modes plus the hypercollisional rate $\mathsf C_{jj}$ of (\ref{tr:coll}). For hyper-resistivity of order $r$, $\eta\,(k_\perp^2/k_{\perp,\max}^2)^r$, and the local estimate $\rho_j\to k_\perp V_j$ for a mode of wavenumber $k_\perp$, condition (\ref{eq:explicitstability}) gives $\Delta t^3\le8\eta/(k_{\perp,\max}V_j)^4$ for $r\le2$ (leaving out the hypercollisional term, which can only relax the condition). For $r=4$ the smallest wavenumber $k_{\perp,\min}$ binds, and the threshold is smaller by the factor $(k_{\perp,\max}/k_{\perp,\min})^{4/3}$. Appendix~\ref{app:timestepping} derives these conditions and tests them through GANDALF's time step. In those tests the slices that become unstable first are the fastest ones, with $|u_j|$ close to $u_{\max}$, and at GANDALF's own time step the IMEX scheme can amplify them by many orders of magnitude while the slow slices change little.

Appendix~\ref{app:timestepping} also establishes two further properties of the IMEX step. GANDALF evaluates the second explicit stage with the Elsasser fields predicted to the midpoint of the step rather than to the stage time. This makes the coupling of $g$ to $z^\pm$, and with it the error of $g$, first order in $\Delta t$. The drift of the invariants remains second order, but it can change sign as $\Delta t$ is varied, so that a coarse scan in $\Delta t$ can suggest a higher order. Among the diagnostics of version 0.6.0, \texttt{hermite\_flux} keeps only $\partial_z$ in the flux (\ref{eq:fluxCD}) and omits the factor $s$ and, at $m=0$, the factor $p_0$ (section~\ref{sec:decorrelation}), and \texttt{hermite\_free\_energy} is not the conserved free energy under the copy closure (section~\ref{tr:formsec}). The omission of the slice speed from the time-step estimate, the first-order coupling of $g$ to $z^\pm$ and the omission of the field-line term from \texttt{hermite\_flux} have been reported in GANDALF's issue tracker.
